\documentclass[11pt,article]{article}

\usepackage[margin=1in]{geometry}
\usepackage[utf8]{inputenc}
\usepackage[T1]{fontenc}
\usepackage{lmodern}
\usepackage{amsmath,amssymb}
\usepackage{hyperref}

\title{\textbf{On Formal Systems and Expressive Rigor in Functional Software}}
\author{Ian Thomas White \\ \texttt{ianthomaswhite\,[at]\,gmail\,[dot]\,com}}
\date{2024}

\begin{document}

\maketitle

\begin{abstract}
This essay explores how type systems in pure functional programming languages serve not merely as verification mechanisms, but as ontological constraints on computational possibility. Drawing on the Curry-Howard correspondence, Martin-Löf type theory, and mathematical structuralism, we argue that the discipline of typed functional programming fundamentally reshapes problem formulation.
\end{abstract}

\section{Constructive Proofs and Executable Specifications}

The Curry-Howard isomorphism establishes a direct bijection between intuitionistic logic and typed lambda calculi: propositions are types, and proofs are terms. Under this framing, writing software is not an act of instructing an abstract machine to mutate register states, but rather the construction of mathematical proofs of well-formedness.

\section{The Ergonomics of Constraint}

In dynamically typed or weakly typed paradigms, safety must be maintained through runtime assertions, defensive programming, and exhaustive integration tests. In contrast, rich type systems (such as Haskell's Hindley-Milner system with GADTs, type families, and higher-kinded types) elevate invariants to the compiler level. By making illegal states unrepresentable, the developer minimizes cognitive load and eliminates entire classes of runtime anomalies.

\section{Conclusion}

Formal systems do not restrict creativity; they provide the scaffolding upon which complex, reliable reasoning can occur without collapse.

\end{document}
